\section{Contexto} \label{sec:contexto}

\subsection{O Dataset PADS e Seu Estudo-Base}

O PADS foi coletado ao longo de três anos no ambulatório de distúrbios do
movimento do Hospital Universitário de Münster, Alemanha
\cite{varghese_machine_2024}. A partir de 504 indivíduos avaliados com um
aplicativo em \textit{smartphone} e dois Apple Watch Series~4 (um em cada punho,
amostragem a 100~Hz em acelerômetro e giroscópio triaxiais), a base pública
consolidou 469 participantes após controle de qualidade e pareamento etário: 276
PD, 79 HC e 114 DD (28 com tremor essencial, 15 parkinsonismo atípico, 11
esclerose múltipla e 60 outros distúrbios) \cite{varghese_machine_2024}.

O protocolo inclui 11 etapas neurológicas padronizadas (oito de 10~s e três de
20~s, $\sim$15~min), gerando 14 janelas de 10~s e 168 canais na formulação
original de ML, além de 30 respostas binárias do PDNMS e metadados clínicos (sem
Hoehn \& Yahr ou MDS-UPDRS na base pública) \cite{varghese_machine_2024}. No
estudo-base~\cite{varghese_machine_2024}, avaliaram-se PD$\times$HC ($n=355$) e
PD$\times$DD ($n=390$) em validação cruzada aninhada 5$\times$5 por participante
(atributos manuais, BOSS, \textit{XceptionTime} e \textit{CatBoost} via
\textit{stacking}), alcançando acurácia balanceada de 91,16\% ($\pm$4,92\%) em
PD$\times$HC e 72,42\% ($\pm$12,33\%) em PD$\times$DD (com verificação pareada por
sexo, $n=234$). O conjunto público encontra-se depositado no repositório
PhysioNet~\cite{PhysioNet-parkinsons-disease-smartwatch-1.0.0}.

\subsection{Revisões Relacionadas}

Revisões recentes sobre IA aplicada à PD (excluídas da síntese primária por CE3)
apontam limitações recorrentes na literatura: um mapeamento de 133 estudos
destacou a baixa padronização entre bases~\cite{shokrpour_machine_2025}; outra
síntese ressaltou a variabilidade entre indivíduos em 11 trabalhos com sensores
vestíveis~\cite{ouabd_overview_2026}; e uma revisão sistemática com 264 artigos
encontrou alto risco de viés em 98,1\% deles e apenas 7,6\% com validação
externa~\cite{pinto_neto_machine_2026}. Ao isolar os \nTotal\ trabalhos que
empregam o PADS, conseguimos comparar recortes metodológicos e riscos de
vazamento sobre exatamente os mesmos dados.
